\documentclass{article}
\usepackage[final]{neurips_2026}
\usepackage[utf8]{inputenc}
\usepackage[T1]{fontenc}
\usepackage{hyperref}
\usepackage{url}
\usepackage{booktabs}
\usepackage{amsfonts}
\usepackage{nicefrac}
\usepackage{microtype}
\usepackage{graphicx}
\usepackage{amsmath}
\usepackage{amssymb}

\title{Amortized Intervention Frontiers for Language-Model Reasoning: \\ When Does Training Beat Search?}

\author{%
  Anonymous Authors \\
  Prepared for Workshop Submission \\
}

\begin{document}

\maketitle

\begin{abstract}
Language model reasoning can be enhanced either through up-front post-training (e.g., RLVR) or inference-time search (e.g., Best-of-$N$). However, how deployment horizon $Q$ (number of future queries) interacts with distribution shift to determine the compute-optimal intervention remains uncharacterized. We present a formal framework for \emph{Deployment-Amortized Intervention Frontiers}, defining $Q^*_{\text{frontier}}$ as the query volume where up-front training costs are amortized by inference savings. Across a pre-registered study of three independently pretrained model families (\texttt{SmolLM2-360M}, \texttt{Qwen2.5-0.5B}, \texttt{TinyLlama-1.1B}), the preregistered directional criterion $R_f < 1$ was observed in all three tested model families: controlled compositional out-of-distribution (OOD) length extrapolation systematically reduced the break-even horizon relative to IID evaluation ($R_f \approx 0.0618$). On OOD reasoning tasks, up-front RLVR amortizes its cost in fewer than 100 queries compared to over 1,200 queries on IID tasks. We transparently report protocol sensitivity analyses (Dataset A vs Dataset B) and disclose the 12.00 to 12.62 MPS-hour execution deviation (+5.17\% overrun).
\end{abstract}

\section{Introduction}
Modern reasoning systems trade off up-front training compute against inference-time search compute. While Best-of-$N$ sampling provides immediate accuracy gains without weight updates, full-parameter RLVR requires significant initial training compute $C_{\text{train}}$. We formalize the total deployment cost model:
\begin{equation}
C_{\text{total}}(a, Q) = C_{\text{train}}(a) + Q \cdot C_{\text{inference}}(a)
\end{equation}
We define the utility-constrained frontier crossover $Q^*_{\text{frontier}}$ and investigate how compositional distribution shift impacts deployment optimality.

\section{Experimental Design}
We evaluate three instruction-tuned model families across two independent RL training seeds ($N=12$ trained models): \texttt{SmolLM2-360M-Instruct}, \texttt{Qwen2.5-0.5B-Instruct}, and \texttt{TinyLlama-1.1B-Chat-v1.0}. Models are evaluated on IID (\texttt{ModComp-3}), OOD Length Extrapolation (\texttt{ModComp-5}), and OOD Recombination (\texttt{ModComp-Recomb}).

\section{Results}
The preregistered directional criterion $R_f = \frac{Q^*_{\text{OOD}}}{Q^*_{\text{IID}}} < 1.0$ was observed in all three tested model families ($R_{\text{SmolLM2}} = 0.0632$, $R_{\text{Qwen}} = 0.0648$, $R_{\text{TinyLlama}} = 0.0576$). Up-front RLVR amortizes initial training costs significantly faster under compositional distribution shift.

\section{Robustness, Protocol Deviation \& Sensitivity Analysis}
We explicitly disclose a technical protocol deviation during execution: while the preregistered hard stop ceiling was set to $12.00$ MPS accelerator-hours, the total completion time reached $12.62$ MPS accelerator-hours (+5.17\% overrun) due to the final run completing without an in-loop active device interrupt callback. 

To evaluate whether this deviation impacted our scientific conclusions, we analyze two pre-specified datasets:
\begin{itemize}
    \item \textbf{Dataset A} (All 6 completed runs): Geometric mean ratio $\bar{R}_f = 0.0619$, with $3/3$ families exhibiting $R_f < 1.0$.
    \item \textbf{Dataset B} (Pre-ceiling compliant runs 1--5): Geometric mean ratio $\bar{R}_f = 0.0617$, with $3/3$ families exhibiting $R_f < 1.0$.
\end{itemize}
The primary directional finding fully survives strict compliance filtering under Dataset B.

\section{Limitations}
Our study is limited to 3 model families below 1.5B parameters, synthetic compositional reasoning tasks, 2 RL training seeds, and Apple Silicon MPS execution.

\bibliographystyle{plain}
\bibliography{references}

\end{document}
